%
% Auteur : Olivier Cots
% Date   : 17 avril 2019
%
% packages latex
%

%---------------------------------------------------------------------------------------------------------------
%\usepackage{wrapfig}
%\renewcommand{\refname}{}

% For notes in the marge
%\usepackage{MnSymbol}
\usepackage{stackengine}

%---------------------------------------------------------------------------------------------------------------
% Voir http://www.cpt.univ-mrs.fr/~masson/latex/Cours-LaTeX-beamer-03.pdf
% sur les packages incontournables pour latexs
\usepackage{import}             % Remplace les inpu et include
\usepackage[french]{babel}      % This package manages culturally determined typographical (and other) rules for a wide range of languages.
%\usepackage[english]{babel}    % This package manages culturally determined typographical (and other) rules for a wide range of languages.
\usepackage[utf8]{inputenc}     % Accept different input encodings
\usepackage[T1]{fontenc}        % Standard package for selecting font encodings
\usepackage{lmodern}            % based on the Computer Modern fonts released into public domain by AMS
\usepackage{verbatim}

%---------------------------------------------------------------------------------------------------------------
% Couleurs
\usepackage{xcolor}             % Driver-independent color extensions for LATEX and pdfLATEX
\usepackage{colortbl}           % The package allows rows and columns to be coloured, and even individual cells.
                                % Utile pour ``Fancy Header Style Options'' dans le main

%---------------------------------------------------------------------------------------------------------------
%
\usepackage{fancyhdr}           % Fancy Header and Footer

\ifSLIDES

    \usepackage{xparse} % pour faire une commande myframe avec 1 ou 0 argument obligatoire

\else

    \usepackage{xparse}
    
    %---------------------------------------------------------------------------------------------------------------
    %---------------------------------------------------------------------------------------------------------------
    % Poly
    
    \usepackage{makeidx}            % Standard package for creating indexe
    %\usepackage{imakeidx} % Pour les index
    
    %---------------------------------------------------------------------------------------------------------------
    % Les chemins ou se trouvent les figures et les extensions de figures
    \usepackage{ifpdf}
    \ifpdf
     \usepackage[pdftex]{graphicx}
     \DeclareGraphicsExtensions{.jpg,.pdf,.png}
     \usepackage[pagebackref,hyperindex=true,linktoc=all]{hyperref}
    \else
     \usepackage{graphicx}
     \DeclareGraphicsExtensions{.ps,.eps}
     \usepackage[dvipdfm,pagebackref,hyperindex=true,linktocpage]{hyperref}
    \fi

    \ifPOLY

        %---------------------------------------------------------------------------------------------------------------
        % Minitoc
        \usepackage[nottoc, notlof, notlot]{tocbibind}
        \usepackage[french]{minitoc}
        %\let\minitocORIG\minitoc
        %\renewcommand{\minitoc}{\minitocORIG \vspace{1.5em}}
        \mtcsettitle{minitoc}{} % pas de titre : Sommaire
        \nomtcrule % pas de trait horizontaux
        %\let\minitocORIG\minitoc
        %\renewcommand{\minitoc}{\minitocORIG \noindent\rule[0.0ex]{1.0\textwidth}{1pt} \vspace{0.0em}}

        %---------------------------------------------------------------------------------------------------------------
        % Les marges
        \usepackage[a4paper,includeheadfoot,margin=2.54cm]{geometry}
        \setlength{\headheight}{14pt}   % ...at least 14.0pt. Pour eviter un warning

    \else

        \ifBOOK

            %---------------------------------------------------------------------------------------------------------------
            % Minitoc
            \usepackage[nottoc, notlof, notlot]{tocbibind}
            \usepackage[french]{minitoc}
            %\let\minitocORIG\minitoc
            %\renewcommand{\minitoc}{\minitocORIG \vspace{1.5em}}
            \mtcsettitle{minitoc}{} % pas de titre : Sommaire
            \nomtcrule % pas de trait horizontaux
            %\let\minitocORIG\minitoc
            %\renewcommand{\minitoc}{\minitocORIG \noindent\rule[0.0ex]{1.0\textwidth}{1pt} \vspace{0.0em}}

            %---------------------------------------------------------------------------------------------------------------
            % Les marges
            %\usepackage[a4paper,includeheadfoot,left=4cm,right=4cm,top=3cm,bottom=3cm]{geometry}
            \usepackage[a4paper, includeheadfoot, margin=2.8cm]{geometry}
            \setlength{\headheight}{14pt}   % ...at least 14.0pt. Pour eviter un warning

        \fi

    \fi

    %---------------------------------------------------------------------------------------------------------------
    % Pour les figures et sous-figures (commande subfigure)
    \usepackage{caption}
    \usepackage{subcaption}

    %---------------------------------------------------------------------------------------------------------------
    % Pour changer le format des titres de section, etc.
    \usepackage[explicit]{titlesec}
    %\usepackage[explicit, newlinetospace]{titlesec}

\fi

%\graphicspath{{\repRessource/figures/}}
\graphicspath{{.}{figures/}{\relativePath/Logos/}}

%---------------------------------------------------------------------------------------------------------------
% Pour pouvoir ajuster les minipages contenant seulement une figure : \includegraphics[width=0.9\textwidth,valign=t]{./bocop_archi}
\usepackage[export]{adjustbox}
%---------------------------------------------------------------------------------------------------------------

%---------------------------------------------------------------------------------------------------------------
% Affichage sous forme de slides : pour reglage de la taille de la page
%\usepackage[screen]{pdfscreen}

%---------------------------------------------------------------------------------------------------------------
% Pour les environnements theoremes, etc.
\usepackage[many, xparse, skins]{tcolorbox}    % http://ctan.org/pkg/tcolorbox

\usepackage{totcount}           % Compteur pour les problèmes
\usepackage{textcomp}           % Commande textasteriskcentered pour newstep : utile dans les preuves pour une separation d'etapes

\@ifpackageloaded{versions}{}{\usepackage{versions}} % Permet de supprimer des environnement a la compilation : excludeversion, ou d'en ajouter : includeversion

\makeatletter
\usepackage{wasysym}            % Pour les pointes noires pleines (RHD) en debut de preuve et autres
\makeatother

%---------------------------------------------------------------------------------------------------------------
% Les maths et ams
\usepackage{bigints}            % big integrale
\usepackage{amsfonts}
\usepackage{amscd}
\usepackage{amsthm}
\usepackage{amsmath}
\usepackage{amssymb}
\usepackage{mathrsfs}           % Donne une alternative a mathcal : mathscr. Les lettres sont encore plus courbees
\usepackage{mathtools}          % Extension package to amsmath. Utile par exemple pour \coloneqq
\usepackage{xstring}            % Pour avoir la commande \IfEqCase, pour la definition des ensembles
\usepackage{dsfont}             % Utile pour la commande \mathds, pour un style different pour la definition des ensembles

\usepackage{stmaryrd}           % Pour les definitions d'intervalles
\usepackage{yhmath}             % L'arc de cercle pour les champs de Jacobi : wideparen

\usepackage[f]{esvect}          % Pour les fleche pour le systeme hamiltonien : \vv{H}

%---------------------------------------------------------------------------------------------------------------
% Les dessins
\usepackage{tikz}
\usepackage{pgfkeys}
\usepackage[european resistor, european voltage, european current]{circuitikz}
\usepackage{pifont}
\usepackage{\relativePathTemplate/tikzgraphicx}
\usetikzlibrary{arrows,shapes,backgrounds,patterns,babel}

%---------------------------------------------------------------------------------------------------------------
% Les tableaux
\usepackage{calc}
\usepackage{array}
\usepackage{booktabs}           % Tableau : epaisseur des lignes


%---------------------------------------------------------------------------------------------------------------
% Les itemize
\usepackage{marvosym}           % Pour la commande Football pour les itemize
\usepackage{enumitem}           % Pour pouvoir redefinir les separation entre items, etc.

%---------------------------------------------------------------------------------------------------------------
% Listings : pour inclure du code matlab, fortran, etc.
\usepackage{listings}

